% GENERATED FILE. Edit canonical lessons, then run npm run generate:books.
% canonical-source-hash: fnv1a64:ec32fb550baccc16
% canonical-lessons: BN-C414-bengk, BN-C414-kauntar, BN-C414-keshiyar, BN-C414-chek, BN-C414-nogod

\chapter{At the Bank}
\label{ch:bn-at-the-bank}
\hlchaptermodality{414}

\begin{chapteropening}
\textbf{By the end of this chapter:} \emph{I can deposit money, queue at a counter and pay by cheque or in cash.}

Complete the last lesson of chapter 414: I can deposit money, queue at a counter and pay by cheque or in cash.
\end{chapteropening}

\section[\texorpdfstring{\bn{ব্যাংক} (bængk)}{ব্যাংক (bængk)}]{\bn{ব্যাংক} (bængk) — a bank}

\label{lesson:BN-C414-bengk}



\begin{quote}
Before the new one: say the Bengali for dirt, then the Bengali for junk.
\end{quote}

\subsection*{You'll want to know: \bn{ব্যাংক}}
\textbf{\bn{ব্যাংক}} — \emph{bængk} — \textquotedblleft{}a bank\textquotedblright{}.

\begin{quote}
\bn{মীরা} \bn{ব্যাংকে} \bn{টাকা} \bn{জমা} \bn{দিল।}
\end{quote}

\begin{quote}
\emph{mīrā bængke ṭākā jômā dilô.}
\end{quote}

\begin{quote}
Mira put money into the bank.
\end{quote}

\begin{cousinweb}
\textbf{\bn{ব্যাংক}} (\emph{bængk}) is borrowed from English \emph{bank}; the word-initial y-phala gives the English vowel æ.
\end{cousinweb}

\begin{grammarlens}[title={What you've built}]
One more word for getting around.
\end{grammarlens}

\subsection*{Your turn}
\begin{itemize}
\raggedright
  \item \emph{Say it:} \emph{bængk}
  \item \emph{Say it:} \emph{bængk}, once more
  \item \emph{Say it:} \emph{mīrā bængke ṭākā jômā dilô.} — Mira put money into the bank.
  \item \emph{Recall:} say the Bengali for dirt, then the Bengali for junk, then say \emph{bængk} again
\end{itemize}

\begin{tcolorbox}[breakable,colback=teal!4,colframe=teal!35!black,title={Before you move on}]
What does \bn{ব্যাংক} mean? (\textquotedblleft{}A bank\textquotedblright{}.) Say it once more.
\end{tcolorbox}
\section[\texorpdfstring{\bn{কাউন্টার} (kāunṭār)}{কাউন্টার (kāunṭār)}]{\bn{কাউন্টার} (kāunṭār) — a counter (in a bank or office)}

\label{lesson:BN-C414-kauntar}



\begin{quote}
Before the new one: say the Bengali for junk, then the Bengali for a bank (for money).
\end{quote}

\subsection*{You'll want to know: \bn{কাউন্টার}}
\textbf{\bn{কাউন্টার}} — \emph{kāunṭār} — \textquotedblleft{}a counter (in a bank or office)\textquotedblright{}.

\begin{quote}
\bn{ব্যাংকের} \bn{দুই} \bn{নম্বর} \bn{কাউন্টারে} \bn{লাইন} \bn{ছিল।}
\end{quote}

\begin{quote}
\emph{bængker dui nômbôr kāunṭāre lāin chhilô.}
\end{quote}

\begin{quote}
There was a queue at counter two in the bank.
\end{quote}

\begin{cousinweb}
\textbf{\bn{কাউন্টার}} (\emph{kāunṭār}) is borrowed from English \emph{counter}.
\end{cousinweb}

\begin{grammarlens}[title={What you've built}]
One more word for getting around.
\end{grammarlens}

\subsection*{Your turn}
\begin{itemize}
\raggedright
  \item \emph{Say it:} \emph{kāunṭār}
  \item \emph{Say it:} \emph{kāunṭār}, once more
  \item \emph{Say it:} \emph{bængker dui nômbôr kāunṭāre lāin chhilô.} — There was a queue at counter two in the bank.
  \item \emph{Recall:} say the Bengali for junk, then the Bengali for a bank (for money), then say \emph{kāunṭār} again
\end{itemize}

\begin{tcolorbox}[breakable,colback=teal!4,colframe=teal!35!black,title={Before you move on}]
What does \bn{কাউন্টার} mean? (\textquotedblleft{}A counter (in a bank or office)\textquotedblright{}.) Say it once more.
\end{tcolorbox}
\section[\texorpdfstring{\bn{ক্যাশিয়ার} (kæshiyār)}{ক্যাশিয়ার (kæshiyār)}]{\bn{ক্যাশিয়ার} (kæshiyār) — a cashier}

\label{lesson:BN-C414-keshiyar}



\begin{quote}
Before the new one: say the Bengali for a bank (for money), then the Bengali for a counter.
\end{quote}

\subsection*{You'll want to know: \bn{ক্যাশিয়ার}}
\textbf{\bn{ক্যাশিয়ার}} — \emph{kæshiyār} — \textquotedblleft{}a cashier\textquotedblright{}.

\begin{quote}
\bn{ক্যাশিয়ার} \bn{আস্তে} \bn{আস্তে} \bn{টাকা} \bn{গুনলেন।}
\end{quote}

\begin{quote}
\emph{kæshiyār āste āste ṭākā gunlen.}
\end{quote}

\begin{quote}
The cashier counted the money slowly.
\end{quote}

\begin{cousinweb}
\textbf{\bn{ক্যাশিয়ার}} (\emph{kæshiyār}) is borrowed from English \emph{cashier}.
\end{cousinweb}

\begin{grammarlens}[title={What you've built}]
One more word for getting around.
\end{grammarlens}

\subsection*{Your turn}
\begin{itemize}
\raggedright
  \item \emph{Say it:} \emph{kæshiyār}
  \item \emph{Say it:} \emph{kæshiyār}, once more
  \item \emph{Say it:} \emph{kæshiyār āste āste ṭākā gunlen.} — The cashier counted the money slowly.
  \item \emph{Recall:} say the Bengali for a bank (for money), then the Bengali for a counter, then say \emph{kæshiyār} again
\end{itemize}

\begin{tcolorbox}[breakable,colback=teal!4,colframe=teal!35!black,title={Before you move on}]
What does \bn{ক্যাশিয়ার} mean? (\textquotedblleft{}A cashier\textquotedblright{}.) Say it once more.
\end{tcolorbox}
\section[\texorpdfstring{\bn{চেক} (chek)}{চেক (chek)}]{\bn{চেক} (chek) — a cheque}

\label{lesson:BN-C414-chek}



\begin{quote}
Before the new one: say the Bengali for a counter, then the Bengali for a cashier.
\end{quote}

\subsection*{You'll want to know: \bn{চেক}}
\textbf{\bn{চেক}} — \emph{chek} — \textquotedblleft{}a cheque\textquotedblright{}.

\begin{quote}
\bn{বাবা} \bn{ভাড়ার} \bn{টাকা} \bn{চেকে} \bn{দিলেন।}
\end{quote}

\begin{quote}
\emph{bābā bhāṛār ṭākā cheke dilen.}
\end{quote}

\begin{quote}
Father paid the rent by cheque.
\end{quote}

\bn{চেকে} \bn{দেওয়া}, to pay by cheque; \bn{চেক} \bn{কাটা}, to write a cheque.

\begin{cousinweb}
\textbf{\bn{চেক}} (\emph{chek}) is borrowed from English \emph{cheque}.
\end{cousinweb}

\begin{grammarlens}[title={What you've built}]
One more word for getting around.
\end{grammarlens}

\subsection*{Your turn}
\begin{itemize}
\raggedright
  \item \emph{Say it:} \emph{chek}
  \item \emph{Say it:} \emph{chek}, once more
  \item \emph{Say it:} \emph{bābā bhāṛār ṭākā cheke dilen.} — Father paid the rent by cheque.
  \item \emph{Recall:} say the Bengali for a counter, then the Bengali for a cashier, then say \emph{chek} again
\end{itemize}

\begin{tcolorbox}[breakable,colback=teal!4,colframe=teal!35!black,title={Before you move on}]
What does \bn{চেক} mean? (\textquotedblleft{}A cheque\textquotedblright{}.) Say it once more.
\end{tcolorbox}
\section[\texorpdfstring{\bn{নগদ} (nôgôd)}{নগদ (nôgôd)}]{\bn{নগদ} (nôgôd) — cash; in cash}

\label{lesson:BN-C414-nogod}



\begin{quote}
Before the new one: say the Bengali for a cashier, then the Bengali for a cheque.
\end{quote}

\subsection*{You'll want to know: \bn{নগদ}}
\textbf{\bn{নগদ}} — \emph{nôgôd} — \textquotedblleft{}cash; in cash\textquotedblright{}.

\begin{quote}
\bn{দোকানদার} \bn{নগদ} \bn{টাকা} \bn{চাইলেন।}
\end{quote}

\begin{quote}
\emph{dokāndār nôgôd ṭākā chāilen.}
\end{quote}

\begin{quote}
The shopkeeper asked for cash.
\end{quote}

Before a noun it means paid in hand: \bn{নগদ} \bn{টাকা}, cash; \bn{নগদে}, in cash.

\begin{cousinweb}
\textbf{\bn{নগদ}} (\emph{nôgôd}) is usually traced to Arabic \emph{naqd}, \textquotedblleft{}ready money\textquotedblright{}, which reached Bengali through Persian.
\end{cousinweb}

\begin{grammarlens}[title={What you've built}]
One more word for getting around.
\end{grammarlens}

\subsection*{Your turn}
\begin{itemize}
\raggedright
  \item \emph{Say it:} \emph{nôgôd}
  \item \emph{Say it:} \emph{nôgôd}, once more
  \item \emph{Say it:} \emph{dokāndār nôgôd ṭākā chāilen.} — The shopkeeper asked for cash.
  \item \emph{Recall:} say the Bengali for a cashier, then the Bengali for a cheque, then say \emph{nôgôd} again
\end{itemize}

\begin{tcolorbox}[breakable,colback=teal!4,colframe=teal!35!black,title={Before you move on}]
What does \bn{নগদ} mean? (\textquotedblleft{}Cash; in cash\textquotedblright{}.) Say it once more.
\end{tcolorbox}
